\chapter[Problema de Pesquisa]{Problema de Pesquisa}

O setor de alimentação fora do lar exige rapidez, controle e organização. Em uma lanchonete, o atendimento ao cliente depende de uma sequência de atividades que começa no recebimento do pedido, passa pela preparação, controle de itens consumidos, eventual atualização do estoque e finaliza no pagamento. Quando essas etapas não são apoiadas por um sistema integrado, a operação se torna mais vulnerável a falhas humanas e a informações desencontradas.

No estabelecimento tomado como referência para este pré-projeto, observa-se a necessidade de organizar melhor processos relacionados a produtos, pedidos, comandas, usuários e relatórios. A execução manual dessas atividades pode provocar erros de anotação, dificuldade para localizar pedidos em andamento, demora no fechamento de contas e baixa visibilidade sobre produtos com maior saída. Além disso, a ausência de um controle formal de acesso dificulta a separação de responsabilidades entre gestor, atendente, caixa e equipe de preparação.

Outro aspecto importante é a confiabilidade das informações. Sem um banco de dados centralizado, registros podem se perder, ser duplicados ou permanecer desatualizados. Essa limitação compromete o acompanhamento do estoque, o planejamento de compras e a análise de vendas. Em pequenos negócios, onde a margem de erro operacional costuma ser reduzida, a falta de informação confiável pode impactar diretamente o atendimento e a sustentabilidade do empreendimento.

Assim, torna-se relevante investigar como uma aplicação web baseada em tecnologias consolidadas pode contribuir para a automação das rotinas internas de uma lanchonete. A escolha de Laravel, PHP, MySQL, Docker, MVC, RBAC e Git/GitHub permite associar recursos de desenvolvimento modernos a uma proposta de baixo custo relativo, adequada ao contexto de pequenos estabelecimentos.

Diante disso, este trabalho busca responder à seguinte questão de pesquisa:

\begin{quote}
\textbf{Como o desenvolvimento de um sistema informatizado utilizando Laravel e MySQL pode contribuir para a automação de processos e para a melhoria da gestão operacional de uma lanchonete?}
\end{quote}